% =============================================================================
% Background
% =============================================================================

\section{Background}
\label{sec:background}

This section establishes the four pieces of vocabulary that
\S\ref{sec:system}--\S\ref{sec:evaluation} assume: vLLM's paged-attention
prefix cache, LMCache's cross-request KV store, cacheblend's selective
recompute, and Claude Code's \texttt{/compact} boundary. The
\texttt{/compact} boundary (\S\ref{sec:bg-claudecode}) is the
integration target.

\subsection{vLLM and prefix caching}
\label{sec:bg-vllm}

vLLM~\cite{vllm-sosp} serves a Transformer LLM by paging its KV cache
into fixed-size blocks (typically 16 tokens). Each block is identified
by a content hash over the token IDs that produced it; when an incoming
request's tokenized prompt has a prefix whose content hashes match
blocks already resident in the cache, vLLM skips the prefill compute
for those blocks and reuses the cached KV directly. Because the hash
is over the exact token sequence and the model weights are fixed,
prefix caching is mathematically lossless: the KV state at the end of
a cache-hit prefill is bit-identical to the state that fresh
recomputation would produce.

This guarantee --- same prefix, same weights, same KV --- is what makes
our no\_cache$\leftrightarrow$prefix\_cache equivalence result in
\S\ref{sec:eval-identity} trivially expected at $T=0$, and what makes
the cross-request mechanism in \S\ref{sec:bg-lmcache} possible.

\subsection{LMCache and cross-request KV reuse}
\label{sec:bg-lmcache}

vLLM's prefix cache lives inside a single engine instance and evicts
under GPU memory pressure. LMCache~\cite{lmcache-paper} extends the
prefix-cache substrate by adding a cross-request, CPU- and
disk-backed store. The lookup semantics are the same (content-addressed
by chunk hash), but the working set is two to three orders of magnitude
larger and persists across requests, sessions, and pod restarts.

We use the V1 connector
(\texttt{LMCacheConnectorV1} in vLLM 0.19+) and its chunk granularity
(256 tokens) throughout. Each chunk that matches the cross-request store
contributes to the request's \emph{cache\_read tokens} counter, which is
the per-turn measurement of \texttt{lmcache}'s engagement with that
request --- we report it directly in \S\ref{sec:eval-hitrate} as the
rescue rate.

\subsection{Cacheblend's selective recompute}
\label{sec:bg-cacheblend}

Plain LMCache requires \emph{exact prefix match}: a single token
inserted at position $k$ invalidates the cache for every chunk from $k$
onwards. This is fine for cache-friendly workloads (RAG with stable
documents, frozen system prompts), but it loses against any workload
that shifts content positionally --- and post-compaction agent traffic
(\S\ref{sec:bg-claudecode}) is exactly that.

Cacheblend~\cite{lmcache-cacheblend} relaxes the exact-prefix
requirement to \emph{positional} match: when a chunk that was the 5th
in one request appears as the 3rd in another, cacheblend reuses the
cached KV for the chunk but recomputes a small fraction of the
attention layers in the rotated positional context (the
\emph{selective recompute}). The relevant knobs:

\begin{itemize}
\item \texttt{blend\_min\_tokens} (256 default): minimum chunk size at
  which the blend mechanism engages. Sub-256-token requests bypass
  cacheblend and fall through to plain prefix caching.
\item \texttt{blend\_recompute\_ratios} (0.15 default): fraction of
  attention layers recomputed per blended chunk. Higher means less
  quality drift but more compute --- the entire point of the mechanism
  is that 0.15 is supposed to be enough.
\item \texttt{blend\_special\_str} (\texttt{\#\,\#} default): the
  string sentinel that marks chunk boundaries in the prompt.
\end{itemize}

The published cacheblend evaluation~\cite{lmcache-cacheblend} reports
70--90\% rescue on permuted-RAG workloads and notes ``small but
bounded'' quality drift, citing accuracy preservation on TruthfulQA at
a single recompute ratio. Agent traffic is not RAG; at $T=0$ greedy
decoding single-token differences are load-bearing; hit rate is not
output quality. \S\ref{sec:evaluation} measures all three.

\subsection{Claude Code and the \texttt{/compact} cliff}
\label{sec:bg-claudecode}

Claude Code~\cite{claude-code} is a coding agent that drives a Claude
model (or, through our proxy, any LLM exposing an Anthropic-shaped
API) through multi-turn tasks. Each turn's request body is, roughly,
the concatenation of: a $\sim$2.5k-token system prompt; the accumulated
conversation history (user messages, model responses, tool calls, and
tool outputs); the user's latest message; and a list of available
tools. As the conversation grows turn over turn, the prompt grows
strictly: turn 5 of a session contains all the tokens of turn 4 plus
the new turn's input and output.

This monotone-growth structure is, by construction, the case prefix
caching is designed for. The pre-compact phase of a CC session
exhibits near-perfect prefix-cache utilization in practice --- each
turn re-prefills only the new tokens of that turn.

The discontinuity is \texttt{/compact}. At a configurable threshold
(default $\sim$80\% of the context window, with a hard floor at 100k
tokens), CC asks the LLM to summarize the conversation so far into an
800-token, 9-section structured summary. The next request then sends:
system prompt + the 800-token summary + a placeholder ``conversation
continues'' marker + the latest user message. The pre-compaction
history is gone; in its place is a synthesized summary that the model
has never seen at this prefix position before.

This is the cliff. From the prefix cache's perspective, the
post-compaction request shares only its system prompt with the
pre-compaction session --- everything after the first $\sim$2.5k
tokens is new. The prefix cache hit ratio collapses to roughly
\texttt{(system\_prompt\_tokens / total\_prompt\_tokens)}, which on
a typical 30k-token post-compact request is around 8\%. The remaining
$\sim$27.5k tokens must be reprefilled, paying a latency hit that
scales linearly with prompt size and a recurring compute cost on every
post-compact turn for the rest of the session.

Cacheblend was designed for this pattern --- summary insertion is
exactly the positional-shift case its selective recompute targets ---
but, as \S\ref{sec:system} shows, adapting cacheblend to real CC
traffic required closing three concrete integration gaps before any
of the rescue numbers in \S\ref{sec:eval-hitrate} held in practice.
